\section{METHODOLOGY}
\label{sec:methodology}

In this section we describe the methodology followed to characterise Galactic foreground emission using multi-frequency angular power spectrum analysis. The approach relies entirely on harmonic-space statistics, avoiding the complications of real-space fitting, and exploits the broad frequency coverage provided by the combination of QUIJOTE, WMAP, and Planck datasets. We first introduce the mathematical framework of angular power spectrum analysis (Section \ref{sec:power_spectrum}), then describe the construction of the analysis masks (Section \ref{sec:masks}), the data processing steps (beam in Section \ref{sec:beam_correction}, unit corrections in Section \ref{sec:unit_conversion}, noise in Section \ref{sec:noise_correction}, and colour corrections in Section \ref{sec:colour_corrections}), the simulations used to estimate uncertainties (Section \ref{sec:simulations}), the foreground models (Section \ref{sec:foreground_model}), and finally, the Bayesian parameter inference (Section \ref{sec:mcmc}).

% ---------------------------------------------------------------
\subsection{Angular power spectrum (APS) analysis}
\label{sec:power_spectrum}
% ---------------------------------------------------------------

The APS is a statistical tool that describes the distribution of signal power as a function of angular scale on the sky. It is a particularly powerful instrument for studying the statistical properties of diffuse emission, since it allows one to separate contributions from different physical processes through their characteristic frequency and multipole dependencies. Rather than working directly in pixel space, all analyses presented in this work are carried out in harmonic space, where the APS constitutes the fundamental observable.

The APS is expressed in terms of the spherical harmonic coefficients $a_{\ell m}$, which encode the sky signal at each multipole $\ell$ and azimuthal order $m$. The multipole $\ell$ is inversely proportional to the angular scale $\theta \sim 180^\circ / \ell$, so that large angular scales correspond to low-$\ell$ values, and small angular scales to high-$\ell$ values. For reference, $\ell = 200$ corresponds to approximately $1^\circ$, which is comparable to the angular resolution of the QUIJOTE-MFI instrument. The cross-power spectrum between two frequency maps $X$ and $Y$ is defined as:

\begin{equation}
    \label{eq:cl_cross}
    C_{\ell}^{X \times Y} = \left\langle \frac{\sum_{m} a_{\ell m}^{X} (a_{\ell m}^{Y})^{*}}{2\ell + 1} \right\rangle,
\end{equation}

where $X$ and $Y$ denote the two frequency bands. When $X = Y$ this reduces to the auto-power spectrum. Cross-spectra probe the signal component common to two independent observations, suppressing the contribution of uncorrelated noise. By combining auto- and cross-spectra from 14 frequency bands (2 QUIJOTE, 5 WMAP and 7 Planck channels), we form a total of 105 spectra (14 auto and 91 cross), greatly improving the constraining power on the foreground model parameters.

All power spectra in this work are computed using the \namaster\ \python\ library\footnote{\url{https://namaster.readthedocs.io/en/latest/}} \citep{Namaster}, which implements the pseudo-$C_\ell$ method with an analytical treatment of the mode-coupling matrix induced by the sky mask. In particular, the E- and B-mode polarization fields are treated with B-mode purification \citep{Smith_2006}, which suppresses E-to-B leakage arising from the application of a partial sky mask. The power spectra are computed in nine multipole bins of width $\Delta\ell = 20$ over the range $20 \leq \ell \leq 200$, with bin edges defined as

\begin{equation}
    \ell_{\rm bin}^{(k)} \in [\ell_1^{(k)},\, \ell_2^{(k)}] = [20 + 20(k-1),\, 39 + 20(k-1)], \quad k = 1, \ldots, 9.
\end{equation}

The lower multipole limit of $\ell = 20$ is imposed by the FDEC filtering applied to the QUIJOTE-MFI maps to remove residual radio frequency interference (RFI) \citep{QUIJOTE_IV}. This filtering suppresses large-angular-scale modes, making multipoles below $\ell \approx 20$ unreliable for a power spectrum analysis. 

The effective multipole $\ell_{\rm eff}$ reported for each bin is computed by \namaster\ as the weighted average of the individual multipoles within the bin:

\begin{equation}
    \ell_{\rm eff}^{(k)} = \frac{\sum_{\ell \in \rm bin\,k} w_\ell\, \ell}{\sum_{\ell \in \rm bin\,k} w_\ell},
\end{equation}

where $w_\ell$ are the per-multipole weights assigned to the bandpower. In this work uniform weights are used, so $\ell_{\rm eff}^{(k)}$ reduces to the arithmetic mean of all integer multipoles in the $k$-th bin, yielding $\ell_{\rm eff} = 29, 49, 69, \ldots, 189$ for the nine bins.


% ---------------------------------------------------------------
\subsection{Masks}
\label{sec:masks}

The analysis is restricted to the sky region accessible to the QUIJOTE-MFI telescope (Section \ref{sec:data_quijote}), which operates from the Teide Observatory and therefore covers a declination band of approximately $-30^\circ < \delta < +90^\circ$. To mitigate contamination from the bright Galactic plane, three analysis masks are constructed by combining the QUIJOTE low-declination survey footprint \citep{QUIJOTE_IV} with Galactic latitude cuts at $|b| > 10^\circ$, $|b| > 15^\circ$, and $|b| > 20^\circ$. Each binary mask is apodized using the C2 scheme of \namaster\ with an apodization scale of $5^\circ$, which tapers the mask smoothly to zero at its edges and suppresses ringing effects in harmonic space. The three masks are shown in Figure~\ref{fig:masks}. The effective sky fractions after apodization are approximately $f_{\rm sky} \approx 0.30$, $0.26$, and $0.23$ for the 10, 15, and 20 degree cuts, respectively. The masks are defined on a \healpix\ grid at $N_{\rm side} = 512$ and are identical for all frequency bands.

\begin{figure}[H]
    \centering
    \begin{minipage}[b]{0.32\linewidth}
        \centering
        \includegraphics[width=\linewidth]{masks/mask_quijote_galcut10_cbar.pdf}
    \end{minipage}
    \hfill
    \begin{minipage}[b]{0.32\linewidth}
        \centering
        \includegraphics[width=\linewidth]{masks/mask_quijote_galcut15_cbar.pdf}
    \end{minipage}
    \hfill
    \begin{minipage}[b]{0.32\linewidth}
        \centering
        \includegraphics[width=\linewidth]{masks/mask_quijote_galcut20_cbar.pdf}
    \end{minipage}

    \caption{Analysis masks combining the QUIJOTE low-declination survey footprint with Galactic latitude cuts of $|b| > 10^{\circ}$ (left), $|b| > 15^{\circ}$ (middle), and $|b| > 20^{\circ}$ (right), after C2 apodisation with a $5^{\circ}$ border.}
    \label{fig:masks}
\end{figure}

% ---------------------------------------------------------------
\subsection{Data corrections}
\label{sec:corrections}
% ---------------------------------------------------------------

The raw power spectra computed from the observed maps contain systematic contributions from instrumental effects that must be accounted for before a physical interpretation can be drawn. In the ideal case, the measured $C_\ell$ coefficients would be proportional to the square of the theoretical spherical harmonic coefficients, $|a_{\ell m}^{\rm theo}|^2$. In practice, several multiplicative and additive corrections are applied, as described below.

\subsubsection{Beam and pixel window functions}
\label{sec:beam_correction}

The finite angular resolution of each instrument and the pixelization of the \healpix\ maps \citep{Gorski_2005} introduce multiplicative window functions that suppress power at high multipoles. The pixel window function, $w_\ell$, arises from the averaging of the sky signal over each pixel and depends solely on $N_{\rm side}$; for $N_{\rm side} = 512$ this correction is small but non-negligible at $\ell > 500$. The instrumental beam is characterised by the window function $B_\ell$, which captures the effective angular resolution of the experiment and causes an exponential suppression of power above the beam scale.

For each frequency band, the beam window function $B_\ell$ is loaded from the official data products of each experiment and interpolated to the effective multipoles of the binning scheme. For QUIJOTE-MFI, the beam is provided as a tabulated array per frequency channel in a FITS file. For WMAP, it is provided as a normalised ASCII table. For Planck LFI and HFI bands, the beam window functions are taken from the respective official FITS products, with separate T, E, and B transfer functions for HFI channels.

\paragraph{}

Both cases, auto- and cross-spectra, are handled by a single correction formula. Defining the effective transfer function for band $i$ as $\mathcal{W}_\ell^i \equiv w_\ell\, B_\ell^i$, where $w_\ell$ is the \healpix\ pixel window function (common to all bands at fixed $N_{\rm side}$) and $B_\ell^i$ is the instrumental beam window function of band $i$, the corrected spectrum between bands $i$ and $j$ is:

\begin{equation}
    \label{eq:corrected_cl}
    C_{\ell}^{i \times j,\,\rm corr} = \frac{C_{\ell}^{i \times j,\,\rm sky} - \delta_{ij}\,N_{\ell}^{i}}{\mathcal{W}_\ell^i\,\mathcal{W}_\ell^j},
\end{equation}

where $\delta_{ij}$ is the Kronecker delta, $N_\ell^i$ is the noise power spectrum of band $i$ (see Section~\ref{sec:noise_correction}), and the denominator deconvolves the combined effect of the pixel window and the beam of each band. The Kronecker delta ensures that the noise bias is subtracted only for auto-spectra ($i = j$); for cross-spectra between different frequency bands the noise contributions of the two detectors are statistically independent and thus average to zero in the cross-correlation, so no noise subtraction is required.


\subsubsection{Unit conversion}
\label{sec:unit_conversion}

All sky maps are converted to a common unit of $\mathrm{K}_{\rm RJ}$ (Rayleigh--Jeans brightness temperature) before computing the power spectra. The unit conversion factor from thermodynamic CMB temperature $\mathrm{K}_{\rm CMB}$ to $\mathrm{K}_{\rm RJ}$ is:

\begin{equation}
    u(\nu) = \frac{c^2}{2 k_B \nu^2} \left.\frac{\partial B_\nu}{\partial T}\right|_{T_{\rm CMB}},
\end{equation}

where $B_\nu$ is the Planck function, $k_B$ is the Boltzmann constant, $c$ the speed of light, and $T_{\rm CMB} = 2.725$~K. For Planck HFI bands (100--353~GHz), the conversion uses the official bandpass-averaged unit conversion coefficients from the Planck 2018 data release \citep{Planck_2018III}, accounting for the instrumental bandpass shape. For all remaining bands the conversion is evaluated analytically at the effective band-centre frequency. The resulting unit conversion factor applied to a power spectrum of the pair $(i,j)$ is $u(\nu_i) \times u(\nu_j)$, where $i = j$ for auto-spectra and $i \neq j$ for cross-spectra.


\subsubsection{Noise power spectrum}
\label{sec:noise_correction}

The instrumental noise contributes an additive bias to the auto-power spectrum that must be estimated and removed (Equation~\ref{eq:corrected_cl}). The noise estimation strategy is specific to each experiment.

For QUIJOTE-MFI, the HMDM is constructed from the two half-survey maps provided in the DR1 data release \citep{QUIJOTE_IV}, together with their associated per-pixel weight maps $w_1$ and $w_2$ (where $w_i = 1/\sigma_i^2$):

\begin{equation}
    \label{eq:hmdm}
    n = \frac{h_1 - h_2}{\sqrt{(w_1 + w_2)(w_1^{-1} + w_2^{-1})}}.
\end{equation}

This construction ensures that the noise power spectrum of the difference map matches that of the combined weighted-sum map, while containing no sky signal \citep{Planck_2015II}.

\paragraph{}

For Planck, the same methodology is followed but the nature of the data split differs between the two focal plane arrays. For the LFI (30-70~GHz), ring-half maps are used as $h_1$ and $h_2$, with pixel covariance maps providing the weights \citep{Planck_2015II}. For the HFI (100-353~GHz), the official half-mission maps from the Planck 2018 data release \citep{Planck_2018III} are used. In both cases, the covariance-weighted HMDM is constructed using Equation~\ref{eq:hmdm}. Since Planck maps are provided at a higher resolution ($N_{\rm side}$), the resulting HMDM is downgraded to $N_{\rm side} = 512$ after construction to match the resolution of the other datasets. 

\paragraph{}

For WMAP, the official band-average maps are noise-weighted co-additions of all Differential Assemblies (DAs) at each frequency. No half-mission split is provided directly for the band-average maps; instead, we construct the HMDM from the temporal split available in the public data release. We use the first four years of observations ($h_1$, years 1 to 4) and the remaining five years ($h_2$, years 5 to 9). Because the two halves average over different numbers of years, the naive difference $h_1 - h_2$ does not reproduce the noise level of the full 9-year combination. We apply the analytical normalization factor that accounts for the unequal time splits:

\begin{equation}
    n_{\rm WMAP} = (h_1 - h_2)\,\sqrt{\frac{t_1\,t_2}{(t_1 + t_2)^2}},
\end{equation}

where $t_1 = 4$ and $t_2 = 5$ are the number of years in each half, giving a normalization factor of $\sqrt{20/81}$. This ensures that the noise power spectrum of the HMDM matches that of the full 9-year map, and it is used directly as $N_\ell^i$ in Equation~\ref{eq:corrected_cl}.

\subsubsection{Colour corrections}
\label{sec:colour_corrections}

Colour corrections account for the mismatch between the SED of the foreground emission and the spectral response of each instrumental bandpass. When a source with spectral index $\beta$ is observed through a band with transmission $\tau(\nu)$, the measured flux density differs from the value at the nominal band-centre frequency by a factor that depends on $\beta$. This factor is the colour correction $\mathrm{CC}(\beta)$.

The colour correction coefficients for each band and each foreground component (synchrotron and thermal dust) are pre-computed using the \fastcc\ code\footnote{\url{https://github.com/mpeel/fastcc}} \citep{Fastcc}, which evaluates the bandpass-weighted ratio between the assumed foreground SED and a flat (Rayleigh-Jeans) reference spectrum.

The treatment differs between foreground components because they follow different SEDs. For synchrotron emission (power law, $I_\nu \propto \nu^\alpha$), the colour correction depends only on a single spectral index $\alpha$, so it can be parameterized as a 1D polynomial. The same applies to dust emission in non-HFI bands, where the modified blackbody curvature is negligible and a power-law approximation is sufficient. For these components, the coefficients are obtained using \fastcc. For thermal dust in the Planck HFI bands (100-353~GHz), however, the SED follows a modified blackbody (MBB),

\begin{equation}
    I_\nu \propto \nu^{\beta_d}\,B_\nu(\nu, T_d),
\end{equation}


where the Planck function $B_\nu(\nu, T_d)$ introduces a non-trivial dependence on the dust temperature $T_d$ that cannot be absorbed into a power-law index. As a result, the HFI dust colour correction depends jointly on both $\beta_d$ and $T_d$, and a simple 1D polynomial in $\alpha$ is insufficient. It is therefore computed using the \interpcc\ interpolator (provided with \fastcc), which interpolates a two-dimensional grid of pre-computed colour corrections as a function of both $\beta_d$ and $T_d$ \citep{Fastcc}. The grid is evaluated at the fixed dust temperature $T_d = 19.6$~K adopted throughout this work, reducing the 2D surface to a 1D function of $\beta_d$ alone, which is then fitted with a second-order polynomial.

\paragraph{}

In all cases the resulting colour correction for component $c \in \{s, d\}$ at frequency band $\nu$ is expressed as a second-order polynomial in the relevant spectral index:

\begin{equation}
    \mathrm{CC}_c(\nu;\, \alpha) = a_0^{c}(\nu) + a_1^{c} (\nu)\, \alpha + a_2^{c}(\nu)\, \alpha^2,
\end{equation}

where $\alpha = \beta_s + 2$ for synchrotron ($c = s$) and $\alpha = \beta_d + 2$ for dust ($c = d$) in Rayleigh-Jeans units. The polynomial coefficients $\{a_0^{c}(\nu),\, a_1^{c}(\nu),\, a_2^{c}(\nu)\}$ are specific to each component $c$ and frequency band $\nu$, and are derived using \fastcc\ following the methodology described above.

\paragraph{}

Rather than being applied as a fixed correction to the data, the colour corrections are folded directly into the foreground model (see Section~\ref{sec:foreground_model}) at each MCMC iteration, using the spectral index proposed by the sampler at that step. Each foreground component carries its own correction, applied at each frequency of the pair. The colour-corrected model terms are:

\begin{equation}
    C_\ell^{\rm sync,\,cc}(\nu_i, \nu_j) = \frac{A_s\,(\ell/\ell_{\rm ref})^{\alpha_s}\,S_s(\nu_i)\,S_s(\nu_j)}{\mathrm{CC}_s(\nu_i;\,\beta_s)\,\mathrm{CC}_s(\nu_j;\,\beta_s)},
\end{equation}

\begin{equation}
    C_\ell^{\rm dust,\,cc}(\nu_i, \nu_j) = \frac{A_d\,(\ell/\ell_{\rm ref})^{\alpha_d}\,S_d(\nu_i)\,S_d(\nu_j)}{\mathrm{CC}_d(\nu_i;\,\beta_d)\,\mathrm{CC}_d(\nu_j;\,\beta_d)},
\end{equation}

and for the synchrotron-dust cross-correlation, each frequency leg receives the correction appropriate for its own component:

\begin{equation}
    C_\ell^{\rm corr,\,cc}(\nu_i, \nu_j) = \rho\sqrt{A_s A_d}\left(\frac{\ell}{\ell_{\rm ref}}\right)^{(\alpha_s+\alpha_d)/2}
    \left[
        \frac{S_s(\nu_i)}{\mathrm{CC}_s(\nu_i;\,\beta_s)}\frac{S_d(\nu_j)}{\mathrm{CC}_d(\nu_j;\,\beta_d)}
        +
        \frac{S_s(\nu_j)}{\mathrm{CC}_s(\nu_j;\,\beta_s)}\frac{S_d(\nu_i)}{\mathrm{CC}_d(\nu_i;\,\beta_d)}
    \right].
\end{equation}

Since the corrections are re-evaluated at every likelihood call, they track the spectral index throughout the chain rather than being fixed at a single approximate value. This propagates the uncertainty in the spectral index self-consistently into the corrected model, improving upon approaches where the colour correction is applied as a fixed factor computed at a mean spectral index \citep{Planck_2018XI}.

% ---------------------------------------------------------------
\subsection{Simulations}
\label{sec:simulations}
% ---------------------------------------------------------------

The statistical uncertainties on the measured power spectra are estimated from Monte Carlo simulations. In the corrected spectrum (Equation~\ref{eq:corrected_cl}), two independent terms contribute to the total error: the uncertainty on the sky power spectrum $C_\ell^{i \times j,\,\rm sky}$, characterized by $\sigma_{C_\ell}$, and the uncertainty on the noise power spectrum $N_\ell^i$, characterized by $\sigma_{N_\ell}$. These are estimated from two independent sets of $N_{\rm sim} = 100$ simulations, and the total uncertainty is obtained as their quadrature sum:

\begin{equation}
    \sigma_\ell = \sqrt{\sigma_{C_\ell}^{2} + \sigma_{N_\ell}^{2}},
\end{equation}

where $\sigma_{C_\ell}$ is the standard deviation of the corrected power spectra across sky-plus-noise realizations, and $\sigma_{N_\ell}$ is the standard deviation of the noise power spectra across noise-only realizations. Note that in the case of cross-spectra, the noise term $N_\ell^i$ does not appear in Equation~\ref{eq:corrected_cl} (since $\delta_{ij} = 0$), so $\sigma_{N_\ell} = 0$ and the total uncertainty reduces to $\sigma_\ell = \sigma_{C_\ell}$ alone.

\subsubsection{Sky signal simulations}
\label{sec:sky_simulations}

Realistic simulated polarization sky maps for each frequency band and experiment are generated using the Python Sky Model (\pysm) package\footnote{\url{https://pysm3.readthedocs.io/en/latest/}} \citep{Pysm}. \pysm\ is a widely used \python\ library for sky emission modelling in CMB analyses (e.g.\ \citealt{Planck_2020LVII, Martire_2022}).

The simulation includes both polarized and intensity foreground components. Since this work focuses exclusively on polarization, only the polarized components enter the power spectrum analysis; the intensity components are included solely to produce realistic total-intensity maps but do not contribute to the polarized signal and are not used in the fitting. The polarized foreground components are: synchrotron emission modelled as a power law with a spatially constant spectral index (\textbf{s1}); thermal dust emission described as a single-component modified blackbody (\textbf{d1}); and a lensed CMB polarization realization (\textbf{c1}), computed using the Taylens code \citep{Taylens} with unlensed $C_\ell$ produced by CAMB\footnote{\url{https://www.camb.info/}}. The intensity foreground components included are spinning dust (AME) modelled as a sum of two spinning dust populations (\textbf{a1}; \citealt{Planck_2015X}), and free-free emission following the analytic Commander model (\textbf{f1}); these are included for completeness but play no role in the polarization analysis.

Simulated sky maps were generated independently for each of the 14 frequency bands using the \pysm\ sky model. For each band, the simulated emission was expressed in thermodynamic CMB temperature units and convolved in harmonic space with the corresponding instrumental effective beam window function. The final maps were produced at a \healpix\ resolution of $N_{\rm side} = 512$. No bandpass integration was performed, each map was generated at the band-centre frequency.

\subsubsection{Noise simulations}
\label{sec:noise_simulations}

Noise simulations are produced independently for each frequency band and experiment, reflecting the distinct noise properties of each instrument.

\subsubsection*{QUIJOTE-MFI}

The noise in the QUIJOTE-MFI maps is spatially correlated and anisotropic, and the two channels within each MFI horn exhibit significant noise correlations both in intensity ($\sim 60$--$80\%$) and in polarization ($\sim 20\%$; \citealt{QUIJOTE_IV}). The QUIJOTE noise simulations used in this work are those described in \citet{QUIJOTE_IV}, which account for the measured anisotropy, spatial correlations, and inter-channel covariance structure.

\subsubsection*{WMAP}

WMAP data are not significantly affected by $1/f$ noise, so the noise is well described by a per-pixel Gaussian model. Following \citet{WMAP_maps}, $N_{\rm sim} = 100$ noise realization maps are generated as independent Gaussian random fields with zero mean and pixel-dependent standard deviation. The per-pixel noise level for polarization is:

\begin{equation}
    \label{eq:wmap_noise}
    \sigma_{\rm pix} = \frac{\sigma_0}{\sqrt{N_{\rm obs,\,eff}}},
\end{equation}

where $\sigma_0$ is the band-level noise sensitivity (tabulated in the WMAP data release) and $N_{\rm obs,\,eff}$ is the effective number of observations per pixel, which accounts for the covariance between the Q and U Stokes parameters stored in the WMAP maps. Specifically, for the Q component:

\begin{equation}
    N_{\rm obs,\,eff}^Q = N_{\rm obs}^Q - \frac{(N_{\rm obs}^{QU})^2}{N_{\rm obs}^U},
\end{equation}

and analogously for U, where $N_{\rm obs}^Q$, $N_{\rm obs}^U$, and $N_{\rm obs}^{QU}$ are the diagonal and off-diagonal elements of the per-pixel observation count matrix given by WMAP for each map and frequency band. This correction ensures that the effective noise per pixel correctly reflects the independent information content of Q and U after accounting for their covariance. The Q and U noise maps are drawn independently using their respective $\sigma_{\rm pix}$.

\subsubsection*{Planck} 

For Planck LFI and HFI bands (PR3 data), $N_{\rm sim} = 100$ per-band noise simulations are downloaded from the Planck Legacy Archive \citep{Planck_2020LVII}. These simulations are generated using the MCMC-based noise estimation procedures described in the respective data processing papers for LFI \citep{Planck_2015II} and HFI \citep{Planck_2018III}, which characterise the noise power spectrum including contributions from the $1/f$ noise component of the detectors. Each simulation map is processed identically to the data in the analysis pipeline.

% ---------------------------------------------------------------
\subsection{Foreground emission model}
\label{sec:foreground_model}
% ---------------------------------------------------------------

The measured power spectra over the frequency range 11-353~GHz are dominated by two Galactic emission components in polarization: synchrotron radiation and thermal dust emission \citep{Planck_2015X}. A cross-correlation term between the two components is included to account for any spatial correlation along the line of sight. This modelling approach is similar to that followed by \citet{Choi_2015} and \citet{Hensley_2022}. The total model power spectrum between frequencies $\nu_i$ and $\nu_j$ we will use in this work is given by Equation \ref{eq:full_model}. 

\begin{equation}
    \label{eq:full_model}
    C_\ell(\nu_i, \nu_j) = C_\ell^{\rm sync}(\nu_i, \nu_j) + C_\ell^{\rm dust}(\nu_i, \nu_j) + C_\ell^{\rm corr}(\nu_i, \nu_j).
\end{equation}

Here $C_\ell^{\rm sync}$ is the synchrotron power spectrum, modelled as a power law in both frequency and multipole space (Section~\ref{sec:synchrotron}). $C_\ell^{\rm dust}$ is the thermal dust power spectrum, modelled via a modified blackbody SED (Section~\ref{sec:dust}); and $C_\ell^{\rm corr}$ is the synchrotron-dust cross-correlation term, which accounts for any spatial correlation between the two components (Section~\ref{sec:corr}).

\subsubsection{Synchrotron emission}
\label{sec:synchrotron}

Synchrotron radiation is the dominant polarized foreground at frequencies below $\sim$70~GHz. Its angular power spectrum is well described by a power law in $\ell$ \citep{Choi_2015, Martire_2022}. Defining the synchrotron frequency scaling as

\begin{equation}
    S_s(\nu) = \left(\frac{\nu}{\nu_{\rm ref}}\right)^{\beta_s},
\end{equation}

\noindent the synchrotron power spectrum is:

\begin{equation}
    \label{eq:synch_model}
C_\ell^{\rm sync}(\nu_i, \nu_j) = A_s \left(\frac{\ell}{\ell_{\rm ref}}\right)^{\alpha_s}
S_s(\nu_i)\, S_s(\nu_j),
\end{equation}

where $A_s$ is the synchrotron amplitude at the reference multipole $\ell_{\rm ref} = 80$ and reference frequency $\nu_{\rm ref} = 23.0$~GHz, $\alpha_s$ is the multipole spectral index, and $\beta_s$ is the frequency spectral index in Rayleigh-Jeans units.

\subsubsection{Thermal dust emission}
\label{sec:dust}

Thermal dust emission dominates polarized foreground emission above $\sim$70~GHz, with a SED following a modified blackbody (MBB) law \citep{Planck_2018XI}. In Rayleigh-Jeans brightness temperature units, the MBB scaling relative to the dust reference frequency $\nu_0 = 353$~GHz is:

\begin{equation}
S_d(\nu) = \left(\frac{\nu}{\nu_0}\right)^{\beta_d} \frac{B_\nu(\nu, T_d)}{B_\nu(\nu_0, T_d)},
\end{equation}

where $B_\nu$ is the Planck function and $T_d = 19.6$~K is the dust temperature, fixed in accordance with Planck measurements \citep{Planck_2015X}. The dust model power spectrum is:

\begin{equation}
    \label{eq:dust_model}
C_\ell^{\rm dust}(\nu_i, \nu_j) = A_d \left(\frac{\ell}{\ell_{\rm ref}}\right)^{\alpha_d}
S_d(\nu_i)\, S_d(\nu_j),
\end{equation}

where $A_d$ is the dust amplitude at $\ell_{\rm ref} = 80$ and $\nu_0 = 353$~GHz, $\alpha_d$ is the angular power spectral index, and $\beta_d$ is the MBB frequency spectral index.

\subsubsection{Synchrotron-dust cross-correlation}
\label{sec:corr}

A non-zero spatial correlation between synchrotron and thermal dust emission can arise when the GMF aligns the polarization angle of both components along the line of sight \citep{Choi_2015}. Following \citet{Hensley_2022}, the cross-correlation term is modelled as:

\begin{equation}
    \label{eq:corr_model}
C_\ell^{\rm corr}(\nu_i, \nu_j) = \rho \sqrt{A_s A_d}\,
\Big[ S_s(\nu_i)\, S_d(\nu_j) + S_s(\nu_j)\, S_d(\nu_i) \Big]
\left(\frac{\ell}{\ell_{\rm ref}}\right)^{(\alpha_s + \alpha_d)/2},
\end{equation}

where $\rho \in [-1, 1]$ is the synchrotron-dust correlation coefficient. When colour corrections are applied, each per-band scaling factor is divided by its corresponding colour correction polynomial, as described in Section~\ref{sec:colour_corrections}.

\subsubsection{CMB subtraction}

Before fitting the foreground model, the CMB contribution is subtracted from the observed power spectra. We use the best-fit theoretical $\Lambda$CDM power spectrum from Planck 2018 \citep{Planck_2018VI}, obtained from the combined Plik TTTEEE likelihood with low-multipole (lowE) and gravitational lensing constraints. Both $C_\ell^{EE}$ and $C_\ell^{BB}$ are subtracted; for BB only the lensing contribution is included, since primordial tensor B-modes have not yet been detected. The CMB spectrum is interpolated to the effective multipoles of the binning scheme before subtraction.

% ---------------------------------------------------------------
\subsection{Bayesian inference and MCMC fitting}
\label{sec:mcmc}
% ---------------------------------------------------------------

We adopt a Bayesian framework to constrain the foreground model parameters from the measured power spectra. The full parameter vector for the power-law fitting is $\boldsymbol{\theta} = \{A_s, \alpha_s, \beta_s, A_d, \alpha_d, \beta_d, \rho\}$, where the EE and BB modes are fitted separately in this work. The posterior distribution of the parameters $\boldsymbol{\theta}$ is

\begin{equation}
p(\boldsymbol{\theta} \mid \mathbf{D}) \propto \mathcal{L}(\mathbf{D} \mid \boldsymbol{\theta})\, p(\boldsymbol{\theta}),
\end{equation}

with $\mathcal{L}(\mathbf{D} \mid \boldsymbol{\theta})$ the likelihood of the data $\mathbf{D}$ given the parameters, and $p(\boldsymbol{\theta})$ the prior. The posterior is thus proportional to the product of the likelihood and the prior, with the normalization constant (the Bayesian evidence) being irrelevant for parameter estimation.



\subsubsection{Likelihood}
\label{sec:likelihood}

The likelihood is constructed assuming that residuals between the measured and model power spectra are Gaussian distributed. This approximation is well justified in the regime of partial sky coverage and moderate multipoles ($\ell < 200$), where the central limit theorem ensures that the distribution of bandpower estimates is close to Gaussian, and has been widely adopted in foreground power spectrum analyses at similar angular scales \citep{Choi_2015, Martire_2022, QUIJOTE_IV}. The data vector $\mathbf{D}$ consists of the measured spectra $\hat{C}_\ell(\nu_i, \nu_j)$ concatenated over all frequency pairs and multipole bins, with corresponding uncertainties $\sigma_\ell(\nu_i, \nu_j)$ estimated from Monte Carlo realizations (Section~\ref{sec:simulations}). The log-likelihood is then:

\begin{equation}
\ln \mathcal{L} = -\frac{1}{2} \sum_{i,j,\ell} \frac{[\hat{C}_\ell(\nu_i, \nu_j) - C_\ell(\nu_i, \nu_j;\boldsymbol{\theta})]^2}{\sigma^2_\ell(\nu_i, \nu_j)}.
\end{equation}

\subsubsection{Priors}
\label{sec:prior}

Two fitting approaches are considered in this work: a power-law fit, in which all nine multipole bins are fitted simultaneously by a global power-law model, and a bin-to-bin fit, in which each multipole bin is fitted independently. Both approaches are described in detail in Section~\ref{sec:fitting_methods}. In both cases, uniform (flat) priors are adopted within physically motivated bounds:

\begin{gather*}
A_s \geq 0, \quad A_d \geq 0, \quad \beta_s \leq 0, \quad \beta_d \geq 0, \quad \rho \in [-1,\, 1],
\end{gather*}

For the bin-to-bin analysis, an additional Gaussian prior on the synchrotron frequency spectral index is included to regularize the fit in multipole bins where the signal-to-noise ratio in the low-frequency channels is too low to adequately constrain the synchrotron SED, which can otherwise lead to spurious results for $\beta_s$ and, consequently, for $\beta_d$ \citep{Planck_2018XI}:

\begin{equation}
\beta_s \sim \mathcal{N}(-3.1,\; 0.30^2),
\end{equation}

consistent with the values found in published WMAP and Planck foreground analyses \citep{WMAP_7yr_foregrounds, Planck_2016XXX}. This approach has also been used in various works as in \cite{Planck_2018XI} or \cite{Choi_2015}.

\subsubsection{Sampler}
\label{sec:sampler}

Posterior sampling is performed with the \emcee\ affine-invariant ensemble sampler \citep{Emcee}. Convergence is assessed for all fits using a combination of standard diagnostics: the integrated autocorrelation time $\hat{\tau}$ is monitored to ensure the post-burn-in chain is sufficiently long; the effective sample size (ESS $= N_{\rm post}/\hat{\tau}$) is required to exceed 1000 per parameter; the mean acceptance fraction across walkers is verified to lie within the range recommended for the affine-invariant sampler; and the reduced chi-squared $\chi_{red}^2$ at the posterior maximum-likelihood point is inspected to confirm that the model provides an adequate fit and that the uncertainties are well calibrated.


% ---------------------------------------------------------------
\subsection{Fitting approaches}
\label{sec:fitting_methods}
% ---------------------------------------------------------------

Two complementary fitting approaches are applied to the corrected power spectra over the multipole range $20 \leq \ell \leq 200$, using data from all 14 frequency bands described in Section~\ref{sec:data}. The baseline analysis includes both auto- and cross-spectra, yielding 105 independent spectra. 

Moreover, a complementary study restricted to cross-spectra only, which avoids any noise bias in the auto-spectra, is presented in Appendix~\ref{sec:appendix_1}; and a study fitting in the multipole range $30 \leq \ell \leq 200$ suggested in \cite{QUIJOTE_IV} is oresented in Appendix~\ref{sec:appendix_2}.

\subsubsection{Power-law fitting}
\label{sec:fitting_methods_pl}

In this approach all nine multipole bins are fitted simultaneously by the global power-law foreground model of Section~\ref{sec:foreground_model}. The MCMC samples the full parameter vector $\boldsymbol{\theta} = \{A_s, \alpha_s, \beta_s, A_d, \alpha_d, \beta_d, \rho\}$ over all frequencies and multipoles. EE and BB modes are fitted separately.

\subsubsection{Bin-to-bin analysis}
\label{sec:fitting_methods_btb}

In this approach each multipole bin is fitted independently, without imposing any global power-law dependence on $\ell$. The per-bin parameter vector is $\boldsymbol{\theta}_{\rm bin} = \{A_s, \beta_s, A_d, \beta_d, \rho\}$. The multipole spectral indices ($\alpha_s$, $\alpha_d$) are not fitted in a single-bin analysis. This approach enables us to examine how the fitted parameters vary with the effective multipole, $\ell_{\rm eff}$, revealing any scale-dependent behavior or local deviations from the global power-law assumption.
